\documentclass[11pt]{amsart}
\usepackage{amsmath,amssymb,amsthm,mathtools}
\usepackage[margin=1in]{geometry}
\usepackage{url}
\usepackage{hyperref}

\newtheorem{theorem}{Theorem}
\newtheorem{lemma}[theorem]{Lemma}
\newtheorem{proposition}[theorem]{Proposition}
\newtheorem{corollary}[theorem]{Corollary}
\theoremstyle{definition}
\newtheorem{conjecture}[theorem]{Conjecture}
\newtheorem{remark}[theorem]{Remark}

\DeclareMathOperator{\vol}{vol}
\newcommand{\Z}{\mathbb{Z}}
\newcommand{\Q}{\mathbb{Q}}
\newcommand{\R}{\mathbb{R}}
\newcommand{\lam}{\lambda}
\newcommand{\hstar}{h^{*}}

\title{Positivity of stretched Littlewood--Richardson\\ coefficients for partitions of length at most four}
\author{Alper Ferudun}
\email{alper@mercurycodelab.com}
\thanks{This is a computer-assisted proof; all finite verifications are
reproduced by the scripts described in Section~\ref{sec:repro}.}
\date{\today}

\begin{document}
\begin{abstract}
For partitions $\lam,\mu,\nu$ the Littlewood--Richardson coefficient
$c^{\nu}_{\lam\mu}$ stretches to a function $P(t)=c^{t\nu}_{t\lam,t\mu}$ which,
by a theorem of Derksen and Weyman, is a polynomial in $t$. King, Tollu and
Toumazet conjectured that $P$ has no negative coefficient. The conjecture is
known only for $c^{\nu}_{\lam\mu}\le 2$ and is otherwise open; it is recorded
as an unsolved problem. We prove it for all triples whose parts number at most
four. The proof is structural: a period-one dilation transfers the
Berline--Vergne local Ehrhart formula from an integral dilate back to the
rational hive polytope, reducing the statement to the positivity of the
Berline--Vergne weight of every two-dimensional transverse cone spanned by
rank-four rhombus normals, which we verify exactly (minimum weight $1/9$).
Vertex integrality and unimodularity of vertex cones are not used, and indeed
fail already in rank four. The same local method yields two rank-uniform
consequences: for every full-dimensional hive polytope of any rank the top four
Ehrhart coefficients are positive, and every full-dimensional five-part hive has
all coefficients positive except possibly the linear one. We close by reducing
the general conjecture to a single hive-specific effectivity statement and
recording the exact obstructions that rule out the standard shortcuts. All
finite verifications are accompanied by independent replay scripts.
\end{abstract}

\maketitle

\section{Introduction}

Let $\lam,\mu,\nu$ be partitions with $|\lam|+|\mu|=|\nu|$ and let
$c^{\nu}_{\lam\mu}$ be the Littlewood--Richardson (LR) coefficient, the
multiplicity of the Schur function $s_\nu$ in $s_\lam s_\mu$. For a positive
integer $t$ write $t\lam=(t\lam_1,t\lam_2,\dots)$. Derksen and Weyman
\cite{DerksenWeyman} proved that
\[
  P^{\nu}_{\lam\mu}(t)\;:=\;c^{t\nu}_{t\lam,t\mu}
\]
agrees, for all $t\ge 1$, with a polynomial in $t$; a short proof was later
given in \cite{ShortPoly}. King, Tollu and Toumazet \cite{KTT2004,KTT2006}
studied these \emph{stretched} coefficients in detail and conjectured a strong
positivity property.

\begin{conjecture}[King--Tollu--Toumazet]\label{conj:ktt}
Every coefficient of $P^{\nu}_{\lam\mu}(t)$, in the monomial basis
$1,t,t^2,\dots$, is nonnegative.
\end{conjecture}

Conjecture~\ref{conj:ktt} is proved for $c^{\nu}_{\lam\mu}=1$
(where $P\equiv 1$, by the Fulton conjecture of Knutson--Tao--Woodward
\cite{KTW}) and for $c^{\nu}_{\lam\mu}=2$ (where $P(t)=t+1$, by Ikenmeyer
\cite{Ikenmeyer} and geometrically by Sherman \cite{Sherman}). Beyond these it
is open, and it is listed as an unsolved problem in the FrontierMath collection
\cite{FrontierMath}. Coquereaux and Zuber \cite{CoqZuber} record the special
case of the linear coefficient in the four-part ($SU(4)$) situation and leave
its nonnegativity as an unproved (``alleged'') assertion. Extensive computation
has produced no counterexample.

We settle the four-part case.

\begin{theorem}\label{thm:main}
Let $\lam,\mu,\nu$ be partitions of length at most four with
$|\lam|+|\mu|=|\nu|$. Then every coefficient of $P^{\nu}_{\lam\mu}(t)$ is
nonnegative. If $c^{\nu}_{\lam\mu}>0$ and the associated hive polytope is
three-dimensional, all four coefficients are strictly positive.
\end{theorem}

The proof is geometric. By the hive theorem of Knutson and Tao \cite{KnutsonTao}
the coefficient $c^{\nu}_{\lam\mu}$ counts the lattice points of the
\emph{hive polytope} $H=H(\lam,\mu,\nu)$, a rational polytope in the space of
interior hive entries; stretching dilates $H$, so $P^{\nu}_{\lam\mu}$ is the
Ehrhart polynomial of $H$ and its degree is $\dim H$. For length at most four,
$\dim H\le 3$. Ehrhart positivity is automatic in dimensions $\le 2$, so the
whole difficulty is the linear coefficient of a three-dimensional hive.

Two features make this delicate. First, hive polytopes for four parts are
lattice polytopes, but this is a nontrivial input; we avoid it entirely by a
dilation argument (Section~\ref{sec:dilation}), which also makes the method work
verbatim in the rational case that occurs for five or more parts. Second,
Ehrhart positivity is \emph{false} for general lattice polytopes---the Reeve
tetrahedra have negative linear coefficient---so any proof must use the special
geometry of hives. Ours uses the Berline--Vergne local Ehrhart formula
\cite{BerlineVergne,RingSchurmann}: the linear coefficient is a sum, over the
edges, of the lattice length of the edge times a weight depending only on the
edge's transverse cone. Because every rank-four rhombus inequality has a
primitive normal with entries in $\{0,\pm1\}$, there are finitely many
transverse cones, and we check that every one of their Berline--Vergne weights
is positive. We emphasise that this is a statement about \emph{edges}: vertex
cones of four-part hives need not be unimodular (Section~\ref{sec:nonsimple}),
but this does not affect an edge-local argument.

The same local computation, carried out uniformly in the number of parts, gives
the following.

\begin{theorem}\label{thm:topfour}
For every full-dimensional hive polytope of intrinsic dimension $D$ (any number
of parts), the coefficients of $t^{D},t^{D-1},t^{D-2},t^{D-3}$ in its Ehrhart
polynomial are positive.
\end{theorem}

\begin{theorem}\label{thm:sidefive}
For every full-dimensional five-part hive polytope (necessarily of dimension
six) all Ehrhart coefficients except possibly the linear one are positive.
\end{theorem}

Theorem~\ref{thm:topfour} is a statement about \emph{full-dimensional} hives;
for such a hive of dimension $D\le 4$ it gives all coefficients, but the only
full-dimensional hives with $D\le 4$ are the four-part ones, so its new content
is the top four coefficients of full-dimensional hives of dimension $D\ge 5$.
Together with Theorem~\ref{thm:main} it leaves, for five parts, only the linear
coefficient of a full-dimensional hive undetermined (Theorem~\ref{thm:sidefive}),
together with the non-full-dimensional five-part hives of intrinsic dimension at
least three, whose face transverse cones are projections of $N_5$ and are not
covered by the atlas.

Finally, in Section~\ref{sec:frontier} we make precise the obstruction to a
general proof. The local formula reduces Conjecture~\ref{conj:ktt} to a
hive-specific ``effectivity'' statement (HTE) about Berline--Vergne/Todd weights
of closed flat-rhombus coarsenings; we prove that termwise (simplex-cone)
positivity is \emph{false} at codimension four, so a general proof must control
the closed coarsenings globally, and we record six standard rank-uniform
shortcuts together with the exact examples that defeat each. These do not
diminish the four-part theorem, which is complete.

\medskip

The finite verifications underlying Theorems~\ref{thm:main}--\ref{thm:sidefive}
are exact (integer and rational arithmetic throughout) and are reproduced by
independent scripts described in Section~\ref{sec:repro}; the Littlewood--%
Richardson counts were cross-checked against two independent tableau/hive
enumerators and against published values \cite{lrcalc}.

\section{The hive model and the local Ehrhart formula}\label{sec:prelim}

We recall the hive description of LR coefficients \cite{KnutsonTao} and fix
coordinates for the four-part case.

A \emph{hive} of size $r$ is a real function on the triangular array
$\{(i,j):0\le j\le i\le r\}$ that is concave across each of the three families
of unit rhombi. With border values
$h(i,0)=\lam_1+\dots+\lam_i$, $h(i,i)=\nu_1+\dots+\nu_i$, and
$h(r,j)=|\lam|+\mu_1+\dots+\mu_j$, the integer hives are in bijection with the
LR tableaux counted by $c^{\nu}_{\lam\mu}$. The interior entries are free
subject to the rhombus inequalities, so the hive polytope
$H(\lam,\mu,\nu)\subset\R^{D}$, $D=\binom{r-1}{2}$, is cut out by a fixed
integer matrix $A_r$ (depending only on $r$) applied to the interior entries,
with a right-hand side that is an integral, homogeneous, linear function of
$(\lam,\mu,\nu)$. Consequently
\[
  H(t\lam,t\mu,t\nu)=t\,H(\lam,\mu,\nu),\qquad
  L_H(t):=\#\bigl(tH\cap\Z^{D}\bigr)=P^{\nu}_{\lam\mu}(t),
\]
so $P^{\nu}_{\lam\mu}$ is the Ehrhart polynomial of $H$ and
$\deg P^{\nu}_{\lam\mu}=\dim H$. For $r=4$ we have $D=3$ and coordinates
$(h_{11},h_{12},h_{21})$; eliminating the border, the $18$ rhombus rows reduce
to $15$ distinct primitive normals, all with entries in $\{0,\pm1\}$
(Section~\ref{sec:rank4}).

We use the Berline--Vergne local Euler--Maclaurin formula
\cite{BerlineVergne} (see also \cite{RingSchurmann} for the presentation we
adopt). For a rational polytope $Q\subset\R^{D}$ and each $k$, the coefficient
of $t^{k}$ in $L_{Q}(t)$ equals
\begin{equation}\label{eq:bv}
  [t^{k}]\,L_{Q}(t)=\sum_{\dim F=k}\alpha\bigl(N(Q,F)\bigr)\,\vol_{\Z}(F),
\end{equation}
where the sum runs over the $k$-dimensional faces $F$, $\vol_{\Z}(F)$ is the
relative lattice volume of $F$, and $\alpha$ is a weight depending only on the
transverse cone $N(Q,F)$ of $Q$ along $F$, and invariant under lattice
translation of that cone. The leading term ($k=D$) is the normalized volume and
the codimension-one term ($k=D-1$) is half the normalized surface area; both are
manifestly positive.

\section{Period-one dilation}\label{sec:dilation}

The next lemma removes any need to know that a hive polytope is a lattice
polytope, and lets us apply the lattice formula \eqref{eq:bv} to the rational
polytopes that occur for five or more parts.

\begin{lemma}\label{lem:dilation}
Fix $r$ and let $N_r$ be the (finite) set of primitive rhombus normals. Suppose
that for every \emph{lattice} polytope $P\subset\R^{D}$ whose facet normals lie
in $N_r$ one has $[t^{k}]L_{P}(t)\ge 0$. Then $[t^{k}]P^{\nu}_{\lam\mu}(t)\ge 0$
for all partitions of length at most $r$.
\end{lemma}

\begin{proof}
Let $H=H(\lam,\mu,\nu)=\{x:A_rx\le b\}$ with $b$ integral, and let $q$ be the
least common multiple of the denominators of the vertices of $H$. Then
$qH=\{x:A_rx\le qb\}$ has the same normal fan as $H$---so its facet normals lie
in $N_r$---and its vertices are $q$ times those of $H$, hence integral; thus $qH$
is a lattice polytope. Since $L_{qH}(t)=\#(t\,qH\cap\Z^{D})=L_H(qt)$, we have
$[t^{k}]L_{qH}=q^{k}\,[t^{k}]L_{H}$ with $q>0$, so the two coefficients have the
same sign. Applying the hypothesis to $qH$ gives $[t^{k}]L_H\ge 0$.
\end{proof}

By Lemma~\ref{lem:dilation} it suffices, for each fixed $r$, to prove Ehrhart
coefficient positivity for lattice polytopes with normals in $N_r$; the
polynomial $P^{\nu}_{\lam\mu}$ inherits it. In particular the four-part theorem
does not depend on the integrality of four-part hives.

\section{The rank-four normal atlas}\label{sec:rank4}

In coordinates $(h_{11},h_{12},h_{21})$ the primitive rank-four rhombus normals
are the $15$ vectors
\[
  \pm e_1,\ \pm e_2,\ \pm e_3,\ \pm(e_1-e_2),\ \pm(e_1-e_3),\ \pm(e_2-e_3),
\]
\[
  (1,-1,-1),\ (-1,1,-1),\ (-1,-1,1),
\]
each with entries in $\{0,\pm1\}$. There are $\binom{15}{2}-6=99$ nonparallel
pairs. Twelve of the fifteen normals are the type-$A_3$ ``alcoved'' directions;
the three remaining ``odd'' rows are what make the atlas fail to be totally
unimodular (over $\binom{15}{3}=455$ triples of rows the determinant multiset is
$\{|{\det}|{=}0{:}146,\ 1{:}272,\ 2{:}36,\ 4{:}1\}$).

\section{Proof of Theorem~\ref{thm:main}}\label{sec:proofmain}

By Lemma~\ref{lem:dilation} we may assume $H$ is a nonempty lattice polytope in
$\R^{3}$ with normals in $N_4$; write $L_H(t)=a_3t^3+a_2t^2+a_1t+1$. (If
$c^{\nu}_{\lam\mu}=0$ then $P\equiv 0$ by saturation \cite{KnutsonTao}, and the
statement is vacuous.) The leading coefficient $a_3$ is the normalized volume
and $a_2$ is half the normalized surface area; both are positive, and if
$\dim H\le 2$ the polynomial is a point, segment, or polygon Ehrhart polynomial,
which is coefficientwise positive. It remains to bound $a_1$ for
$\dim H=3$.

Apply \eqref{eq:bv} with $k=1$: since a three-dimensional polytope has each edge
$E$ as a ridge lying in exactly two facets, the transverse cone $N(H,E)$ is the
two-dimensional cone spanned by the two facet normals meeting along $E$, both in
$N_4$. Hence
\[
  a_1=\sum_{\text{edges }E}\alpha\bigl(N(H,E)\bigr)\,\ell_{\Z}(E),
\]
with $\ell_{\Z}(E)\ge 0$ the lattice length of $E$.

\begin{proposition}\label{prop:codim2}
For every ordered pair of nonparallel normals in $N_4$, the Berline--Vergne
weight $\alpha$ of the corresponding pointed two-dimensional transverse cone is
positive; over the $99$ pairs its minimum value is $1/9$.
\end{proposition}

Proposition~\ref{prop:codim2} is a finite exact computation
(Section~\ref{sec:repro}); the weight of a saturated normal basis with Gram
matrix $\left(\begin{smallmatrix}A&C\\ C&B\end{smallmatrix}\right)$ equals
$\tfrac14-\tfrac{C}{12}\bigl(\tfrac1A+\tfrac1B\bigr)$, and evaluating this over
the atlas gives values in $\{1/9,\,5/18,\,1/4\}$, all positive. Granting it,
\[
  a_1\ \ge\ \tfrac19\sum_{\text{edges }E}\ell_{\Z}(E)\ >\ 0,
\]
so all four coefficients of $L_H$ are positive. Undoing the dilation
(Lemma~\ref{lem:dilation}) proves Theorem~\ref{thm:main}. \qed

\subsection{Nonsimple vertices do not obstruct the argument}\label{sec:nonsimple}
The hypotheses use only that $a_1$ is edge-local. For definiteness we note that
four-part hives really do have nonunimodular vertex cones: at
$\lam=\mu=(12,8,4)$, $\nu=(18,14,10,6)$ the vertex $(26,32,38)$ is simple with
primitive tangent rays $(0,1,1),(1,0,1),(1,1,0)$, whose determinant is $2$
(and the polytope has $c^{\nu}_{\lam\mu}=50$, all vertices integral,
$L_H(t)=1+7t+18t^2+24t^3$). This has no bearing on \eqref{eq:bv} for $k=1$: each
edge still lies in exactly two facets, and any further rhombus rows tight along
an edge are redundant on, or interior to, its two-dimensional transverse cone.

\section{Rank-uniform coefficients and five parts}\label{sec:uniform}

Formula \eqref{eq:bv} localizes each coefficient on faces of a fixed codimension,
whose transverse cones, for hives of any rank, are generated by the fixed normal
set $N_r$. Enumerating the possible closed transverse cones and evaluating
$\alpha$ exactly gives Theorem~\ref{thm:topfour}: the volume and surface terms
handle codimensions $0,1$, and exact rank-uniform positivity of $\alpha$ on every
closed codimension-two and codimension-three hive transverse cone handles $t^{D-2}$
and $t^{D-3}$.

For Theorem~\ref{thm:sidefive} ($r=5$, $D=6$) one further coefficient, the
quadratic, is controlled at codimension four. Here termwise positivity already
fails---an individual saturated four-normal cell can have weight $-349/28800$---%
but the exact slack closure forces such a cell either to raise the normal rank to
at least five or to embed in one of finitely many pointed rank-preserving closed
supersets, and every one of those full cones has positive weight (minimum
$739/86400$). Enumerating all $\binom{27}{4}=17550$ four-normal tuples and their
closures (903 nonsaturated tuples with minimum weight $1/14400$; 132 negative
saturated cells, each repaired) yields positivity of $t^{6},\dots,t^{2}$, leaving
only the linear coefficient. The full enumeration and its independent replay are
in the codimension-four report.

Theorem~\ref{thm:sidefive} concerns full-dimensional five-part hives. The
lower-dimensional five-part hives are not reduced to four parts by Horn-facet
factorization: the triple $\lam=(4,3,3,1)$, $\mu=(4,2,1,1)$, $\nu=(6,5,4,2,2)$
has $\dim H=3$ while all $142$ essential Horn inequalities are strict.

\section{The general frontier}\label{sec:frontier}

For a period-one hive polytope the local formula writes every Ehrhart coefficient
as a pairing of nonnegative face volumes with Berline--Vergne/Todd weights on
transverse cones. Thus Conjecture~\ref{conj:ktt} follows from a single
hive-specific effectivity statement: for every rank, codimension $q$, and
$2$-connected primitive-interior closed flat-rhombus coarsening $\Sigma$ with
weight vector $a_q$ and Minkowski boundary map $\partial_q$,
\begin{equation}\label{eq:hte}
  \exists\, y:\quad a_q+\partial_q^{\mathsf T}y\ \ge\ 0. \tag{HTE}
\end{equation}
By rational Farkas duality \eqref{eq:hte} is equivalent to
$\langle a_q,w\rangle\ge 0$ for every nonnegative balanced realizable face weight
$w$; pairing it with the (positive) face-volume weight of any hive polytope makes
every coefficient nonnegative. Horn factorization \cite{KTT2006,KTTfactor}
disposes of separating boundary coarsenings, so the load-bearing case is the
$2$-connected primitive interior. We have neither a proof of \eqref{eq:hte} nor a
negative realizable balanced weight refuting it.

On the disproof side, one transfer route deserves record because it is
\emph{valid} and still fails. Realizing three disjoint horizontal strips as a
skew shape identifies the lattice points of any $3\times N$ transportation
polytope with Littlewood--Richardson tableaux of a fixed triple, compatibly
with dilation: for example, the polytope with row margins $(3,3,3)$ and column
margins $(2,1^{7})$ has $L_T(n)=c^{\,n\nu}_{n\lam,\,n\mu}$ for
$\lam=(9,6,3)$, $\mu=(9,9,7,6,5,4,3,2,1)$, $\nu=(15,12,9,7,6,5,4,3,2,1)$
(both counts equal $1050$ at $n=1$). Hence a $3\times 8$ transportation
polytope with a negative Ehrhart coefficient would refute
Conjecture~\ref{conj:ktt} outright. None exists in the relevant regime: over
the full dimension-$14$, codegree-$3$ margin family the minimum linear
coefficient is $2157/280>0$ (exhaustive for total margin $\le 19$, with
saturation in every unbounded direction, and with the closed form
$\tfrac{3a}{2}\bigl(1+H_{2a-1}-H_a\bigr)>0$ for the symmetric unit-column
subfamily); moreover codegree $3$ forces $L_T(1)\ge 1050$, so the known
negative order polytope of Liu and Tsuchiya \cite{LiuTsuchiya}, which has
$255$ lattice points, is not realizable in this shape. The route evades the
positivity theorems for skew Gelfand--Tsetlin polytopes (its content
constraints are global sums, not markings) and dies only on these measured
facts.

Termwise (simplex-cone) positivity is genuinely false---the $-349/28800$ cell
above---so a general proof cannot proceed cell by cell; the finite success at
codimension four does not extend to an unbounded certificate cascade. We record,
finally, that the standard rank-uniform shortcuts are individually obstructed by
exact examples: generic type-$A$ Todd/cocircuit balancing (a nonnegative
$A_4$ Dahmen--Micchelli polynomial maps to $N^4+4N^3+3N^2-2N-\tfrac65$);
network-flow equivalence (simple hive tangent cones have determinant two);
generic Cohen--Macaulay/Koszul Hilbert-ring structure (the $Q_{20}$ Hibi ring is
Gorenstein, Koszul, with negative $a$-invariant, yet its Hilbert polynomial has
linear coefficient $-168011/330$); matroidal/secondary-fan Todd effectivity
(a genuine hive face carries a $U_{2,4}$ tight-normal restriction); positive
tableau/Kostka recursions (they require signs, or stop on primitive triples);
and rank-by-rank Hurwitz stability (an unbounded hierarchy). Details and the
supporting computations are in the accompanying reports.

\section{Reproducibility}\label{sec:repro}

Every finite claim above is verified by exact arithmetic and reproduced by
scripts distributed with this paper.
\begin{itemize}
\item Two independent Littlewood--Richardson counters (a hive lattice-point
enumerator and a lattice-word LR-tableau counter), cross-checked against each
other on all triples with $|\nu|\le 8$, against $c=1,2$ stretched values, and
against \cite{lrcalc}.
\item The rank-four atlas of $15$ normals and $99$ pairs, and
Proposition~\ref{prop:codim2} (all Berline--Vergne weights, minimum $1/9$).
\item An independent Ehrhart interpolation that recomputes $L_H(0),\dots,L_H(D+2)$
and confirms the interpolated polynomial at two held-out points, checked against
both LR counters; and a verification of $a_1=\sum_E\alpha\,\ell_{\Z}(E)$ against
lattice-count $a_1$ on more than $10^4$ four-part hives that are not among the
atlas basis polytopes.
\item The rank-uniform codimension-two and codimension-three enumerations, and
the full codimension-four five-part enumeration with the closure repairs.
\end{itemize}
The dilation identity of Lemma~\ref{lem:dilation} and the nonunimodular vertex of
Section~\ref{sec:nonsimple} are likewise reproduced exactly.

The complete reproducibility bundle---the two Littlewood--Richardson counters,
the rank-four atlas and codimension-two/three enumerations, the five-part
codimension-four certificate and its independent checker, and the exact replay
scripts for every displayed value---is available at
\url{https://github.com/AlperTheKing/ktt-positivity}.

\begin{thebibliography}{99}
\bibitem{BerlineVergne} N.~Berline and M.~Vergne, \emph{Local Euler--Maclaurin
formula for polytopes}, Moscow Math.\ J.\ \textbf{7} (2007), 355--386,
\href{https://arxiv.org/abs/math/0507256}{arXiv:math/0507256}.
\bibitem{CoqZuber} R.~Coquereaux and J.-B.~Zuber, \emph{From orbital measures to
Littlewood--Richardson coefficients and hive polytopes}, Ann.\ Inst.\ Henri
Poincar\'e D \textbf{5} (2018), 339--386,
\href{https://arxiv.org/abs/1706.02793}{arXiv:1706.02793}.
\bibitem{DerksenWeyman} H.~Derksen and J.~Weyman, \emph{On the
Littlewood--Richardson polynomials}, J.\ Algebra \textbf{255} (2002), 247--257.
\bibitem{FrontierMath} Epoch AI, \emph{FrontierMath open problems: stretched
Littlewood--Richardson coefficients},
\href{https://epoch.ai/frontiermath/open-problems/stretched-lr-coefficients}{epoch.ai}
(accessed 2026-07-22).
\bibitem{Ikenmeyer} C.~Ikenmeyer, \emph{Small Littlewood--Richardson
coefficients}, J.\ Algebraic Combin.\ \textbf{44} (2016), 1--29.
\bibitem{KnutsonTao} A.~Knutson and T.~Tao, \emph{The honeycomb model of
$GL_n(\mathbb{C})$ tensor products I: proof of the saturation conjecture},
J.\ Amer.\ Math.\ Soc.\ \textbf{12} (1999), 1055--1090.
\bibitem{KTW} A.~Knutson, T.~Tao and C.~Woodward, \emph{The honeycomb model II:
Facets of the Littlewood--Richardson cone}, J.\ Amer.\ Math.\ Soc.\ \textbf{17}
(2004), 19--48.
\bibitem{KTT2004} R.~C.~King, C.~Tollu and F.~Toumazet, \emph{Stretched
Littlewood--Richardson and Kostka coefficients}, in \emph{Symmetry in Physics},
CRM Proc.\ Lecture Notes \textbf{34} (2004), 99--112.
\bibitem{KTT2006} R.~C.~King, C.~Tollu and F.~Toumazet, \emph{The hive model and
the polynomial nature of stretched Littlewood--Richardson coefficients},
S\'em.\ Lothar.\ Combin.\ \textbf{54A} (2006), B54Ad.
\bibitem{KTTfactor} R.~C.~King, C.~Tollu and F.~Toumazet, \emph{Factorisation of
Littlewood--Richardson coefficients}, J.\ Combin.\ Theory Ser.\ A \textbf{116}
(2009), 314--333.
\bibitem{LiuTsuchiya} F.~Liu and A.~Tsuchiya, \emph{Stanley's non-Ehrhart-positive
order polytopes}, Adv.\ in Appl.\ Math.\ \textbf{108} (2019), 1--10,
\href{https://arxiv.org/abs/1806.08403}{arXiv:1806.08403}.
\bibitem{lrcalc} P.~Alexandersson, \texttt{lrcalc-rs},
\href{https://github.com/PerAlexandersson/lrcalc-rs}{github.com/PerAlexandersson/lrcalc-rs}.
\bibitem{RingSchurmann} M.~H.~Ring and A.~Sch\"urmann, \emph{Local formulas for
Ehrhart coefficients from lattice tiles}, Beitr.\ Algebra Geom.\ \textbf{61}
(2020), 157--185, \href{https://arxiv.org/abs/1709.10390}{arXiv:1709.10390}.
\bibitem{Sherman} C.~Sherman, \emph{Geometric proof of a conjecture of King,
Tollu and Toumazet}, \href{https://arxiv.org/abs/1505.06551}{arXiv:1505.06551}.
\bibitem{ShortPoly} \emph{A short proof of the polynomiality of stretched
Littlewood--Richardson coefficients},
\href{https://arxiv.org/abs/2211.06810}{arXiv:2211.06810}.
\end{thebibliography}

\end{document}
